% SOURCE_AUTHORITY: sga5_fr_workpass.tex, lines 15302--15359 (LF-normalized unit).
% SOURCE_SHA256: 791F4EFFC5E02832D5D77ED03518C8156D6F07E4C8238B03545DB93D883FBB28
\section*{Índice terminológico}
\newcommand{\sgaindexentry}[2]{\noindent\begin{minipage}[t]{0.62\linewidth}#1\end{minipage}\hfill\begin{minipage}[t]{0.32\linewidth}#2\end{minipage}\par\smallskip}
\sgaindexentry{bidualidade (teorema de)}{I 3.4.1}
\sgaindexentry{$A$-categoria, $A$-categoria abeliana}{V 1.1}
\sgaindexentry{classe de $(c,N)$, classe de $c$}{III 3.2.4}
\sgaindexentry{classe de cohomologia associada a $Y$}{III 6.8}
\sgaindexentry{classe de Chern}{VII 3.2, 3.9}
\sgaindexentry{composição (de correspondências) }{III 5.2}
\sgaindexentry{Condição (reg) }{I 3.1.1}
\sgaindexentry{condição de Mittag-Leffler-Artin-Rees}{V 2.1.1}
\sgaindexentry{conjectura de divisibilidade de Grothendieck}{III B 6, XII 4.5}
\sgaindexentry{feixe constante torcido ($\ZZ_\ell$-feixe)}{VI 1.2.1}
\sgaindexentry{correspondência cohomológica}{III 3.1, 3.2, 6.4}
\sgaindexentry{correspondência equivariante}{III B 6}
\sgaindexentry{correspondência diagonal}{III 3.2 b)}
\sgaindexentry{correspondência de Frobenius}{III 3.2 b), XV 2.1}
\sgaindexentry{produto cup de correspondências}{III 4.1.4}
\sgaindexentry{dimensão quase-injetível}{I 1.4}
\sgaindexentry{divisor com cruzamentos normais}{I 3.1.5}
\sgaindexentry{dual (correspondência)}{III 5.1}
\sgaindexentry{dualizante (complexo)}{I 1.7}
\sgaindexentry{dualidade local}{I 4.2}
\sgaindexentry{Euler--Poincaré (característica de) }{VII 4.11}
\sgaindexentry{Euler--Poincaré (fórmula de) }{X 7.1}
\sgaindexentry{feixe $\ell$-ádico construtível}{VI 1.1.1}
\sgaindexentry{$\ZZ_\ell$-feixe}{VI 1.1.4}
\sgaindexentry{$A$-feixe construtível}{VI 1.4.1, 1.4.3}
\sgaindexentry{$\QQ_\ell$-feixe construtível, constante torcido}{VI 1.4.2}
\sgaindexentry{$AR$-$\ZZ_\ell$-feixe construtível}{VI 1.5.1}
\sgaindexentry{flecha de tipo $i$ sobre $f$}{III 1.4}
\sgaindexentry{flecha de Künneth}{III 1.6, 1.7}
\sgaindexentry{flecha de avaliação}{III 2.2.2, III B 5.3.5}
\sgaindexentry{flecha de produto tensorial}{III B 5.4.2}
\sgaindexentry{fortemente dessingularizável}{I 3.1.5}
\sgaindexentry{Frobenius (endomorfismo de)}{XV 1}
\sgaindexentry{Frobenius (correspondência de)}{III 3.2 b), XV 2.1}
\sgaindexentry{Frobenius relativo (morfismo de)}{XV 2 def. 3}
\sgaindexentry{Frobenius aritmético, geométrico}{XV p. 20}
\sgaindexentry{Gauss--Bonnet (fórmula de) }{VII 4.9}
\sgaindexentry{Gysin (sequência exata de)}{VII 1.5}
\sgaindexentry{imagem direta (de correspondência cohomológica) }{III 3.3}
\sgaindexentry{Lefschetz--Verdier (fórmula de) }{III 4.5.1, 4.7, 6.7.2}
\sgaindexentry{Nielsen--Wecken (invariante local de)}{XII 3.1}
\sgaindexentry{Nielsen--Wecken (fórmula de) }{XII 5.1}
\sgaindexentry{normalizado em $x$ (complexo dualizante)}{I 4.1, 4.5}
\sgaindexentry{produto tensorial externo (de correspondências)}{III 5.3}
\sgaindexentry{pureza cohomológica absoluta}{I 3.1.4, I Ap. 4.4}
\sgaindexentry{quase-injetivo}{I 1.1}
\sgaindexentry{racionalidade das funções $L$ (teorema de) }{XV 2 n\textsuperscript{o} 2}
\sgaindexentry{representação de Artin}{X 4.2.5}
\sgaindexentry{representação de Swan}{X 4.2.2}
\sgaindexentry{autointerseção (fórmula de) }{VII 4.1, 9.2}
\sgaindexentry{Shih (teorema de)}{V A 3.1}
\sgaindexentry{sistema projetivo AR-nulo}{V 2.2.1}
\sgaindexentry{sistema projetivo $J$-ádico, AR-$J$-ádico}{V 3.1.1, 3.2.2, 5.1.1, 5.1.3}
\sgaindexentry{termo local de Lefschetz--Verdier}{III 4.2.7, 6.6.3, III B 6}
\sgaindexentry{traço}{III B 5.6.1, 6}
\sgaindexentry{Woods Hole (fórmula de) }{III 6.12}
